\documentclass[11pt,a4paper,twoside,openright]{book}

% Including necessary packages
\usepackage[a4paper,left=2.5cm,right=2.5cm,top=2.5cm,bottom=2.5cm]{geometry}
\renewcommand{\baselinestretch}{1.15}
\usepackage{hyperref}
\usepackage{graphicx}
\usepackage{caption}
\usepackage{fontspec}
\usepackage{fancyhdr}
\usepackage[nottoc]{tocbibind}
\usepackage[toc,page]{appendix}
\usepackage[backend=biber,style=authoryear,maxcitenames=1,autocite=inline,natbib=true]{biblatex}
\addbibresource{tese.bib}
\usepackage{indentfirst}
\usepackage{pdfpages}
\usepackage[portuguese,english]{babel}
\usepackage{filecontents}
\usepackage{ragged2e}

% Configuring headers and footers
\pagestyle{fancy}
\fancyhf{}
\fancyhead[RO]{\nouppercase{\rightmark}}
\fancyhead[LE]{\nouppercase{\leftmark}}
\fancyfoot[C]{\thepage}
\renewcommand{\headrulewidth}{0pt}
\renewcommand{\footrulewidth}{0pt}

% Setting font for XeLaTeX
\setmainfont{Times New Roman}

% Defining custom commands for keywords
\providecommand{\keywords}[1]{\textbf{Keywords:} #1}
\providecommand{\keywordsP}[1]{\textbf{Palavras-chave:} #1}

\begin{document}

\frontmatter

% Adding the cover page as a PDF
% \begin{titlepage}
% \includepdf{capa.pdf}
% \end{titlepage}

% Including abstracts
\chapter*{Abstract}
\selectlanguage{english}
A concise summary (250-300 words) of the thesis, introducing the problem of predicting judicial decisions and detecting errors in Portuguese law. The methodology includes LLMs, RAG, CoT, XAI, fine-tuning, and lawyer validation. Contributions advance legal tech and judicial efficiency in Portugal, with significance in academic novelty and practical applications.
\keywords{LLMs, judicial prediction, error detection, legal tech, Portuguese law}

\chapter*{Resumo}
\selectlanguage{portuguese}
Um resumo conciso (250-300 palavras) da tese, apresentando o problema de prever decisões judiciais e detetar erros na lei portuguesa. A metodologia inclui LLMs, RAG, CoT, XAI, ajuste fino e validação por advogados. As contribuições avançam a tecnologia jurídica e a eficiência judicial em Portugal, com significado em novidade acadêmica e aplicações práticas.
\keywordsP{LLMs, predição judicial, deteção de erros, tecnologia jurídica, direito português}

% Adding table of contents, figures, and tables
\selectlanguage{english}
\renewcommand{\contentsname}{Table of Contents}
\tableofcontents
\listoffigures
\listoftables

\mainmatter

% Chapter 1: Introduction
\chapter{Introduction}
\section{Background and Motivation}
Explores the global rise of AI in legal systems and addresses challenges in the Portuguese judicial context (e.g., language, civil law system).
\section{Problem Statement}
Defines the dual focus: predicting decisions and detecting errors in judicial rulings, specifying domains: penal law, child and family law, labor law (Portuguese context).
\section{Research Questions, Hypotheses, and Objectives}
\subsection{Research Questions}
E.g., "Can LLMs accurately predict outcomes in Portuguese penal law cases?" "How effectively can LLMs detect judicial errors in child and family law?"
\subsection{Hypotheses}
E.g., "Fine-tuned LLMs with RAG outperform traditional models in prediction accuracy." "CoT prompting enhances error detection interpretability."
\subsection{Objectives}
Academic: Develop a robust model for prediction and error detection. Practical: Explore applications for improving judicial efficiency and transparency.
\section{Scope and Limitations}
Focuses on penal, child and family, and labor law cases from dgsi.pt. Limitations include data availability, GPU constraints (8GB VRAM), and limited law firm availability.
\section{Thesis Structure}
Provides an overview of each chapter’s purpose and content.

% Chapter 2: Literature Review
\chapter{Literature Review}
\section{Introduction to Literature Review}
Defines scope: LLMs, legal tech, and Portuguese law.



%%%%%%%%%%%%%%%%%%%%%%%%%%%%%%%%%%%%%%%%%%%%%%%%%%%%

\section{Portuguese Legal System}
% Provides an overview of penal, child and family, and labor law structures. Alongside the "judicial error" legal paradigm in portugal.

Portugal is a Democratic State of Law, where every citizen is assured the right of access to the law and to the courts of justice for defense of their rights and interests.

\par Tribunals are the sovereign entities to administer justice in the name of people and are organized as follows by the Portuguese Constitution and Law 62/2013 of August 26th \parencite{noauthor_lei_nodate,noauthor_dgpj_nodate}:

\begin{itemize}
    \item \textbf{Constitutional Court:} assigned to administer justice in matters of constitutionality, enforcing the supremacy of the Portuguese Constitution.
    \item \textbf{Supreme Court of Justice:} highest court for judicial matters, with jurisdiction over all other judicial courts in all the Portuguese territory.
    \item   \textbf{Courts of Appeal:} these tribunals are the second instance, where appeals from the first instance courts are judged. They are divided into 5 regions: Lisbon, Porto, Guimarães, Coimbra and Évora.
    \item \textbf{Courts of First Instance:} these are the courts where cases initialize . They are divided into several competences, depending on the nature of the case:
            \begin{itemize}
                \item \textbf{Civil Courts:} handle civil matters such as contracts, property disputes, family law, etc.
                \item \textbf{Criminal Courts:} deal with criminal cases, including felonies and misdemeanors.
                \item \textbf{Family and Minors Courts:} administer cases related to family law, child protection, custody, and adoption.
                \item \textbf{Labor Courts:} specialize in labor disputes between employers and employees.
                \item \textbf{Commercial Courts:} focus on commercial and business-related disputes.
                \item \textbf{Execution Courts:} responsible for fiscalizing the execution of judicial decisions. This includes enforcing or extinguishing prison/expulsion sentences, emmiting arrest or liberation warrants, among other powers.
                \item \textbf{Specialized Courts:} such as Maritime Courts, Intellectual Property Courts, among other specialized jurisdictions.
            \end{itemize}
    \end{itemize}

Administrative and Fiscal Courts have a similar hierarchical structure to the Judicial Courts:

\begin{itemize}

    \item \textbf{Supreme Administrative Court:} highest court for administrative and fiscal matters.
    \item \textbf{Central Administrative Courts:} these tribunals are the second instances for administrative and fiscal cases.
    \item \textbf{Administrative and Fiscal Courts of First Instance:} these are the courts where administrative and fiscal cases initialize.        
\end{itemize}

Additionally, there are other tribunals with focus on specific matters like:
\begin{itemize}
    \item \textbf{Courts of Auditors:} tasked with fiscalizing public expenditure.
    \item \textbf{Military Courts:} used only in times of war, however, in times of peace there are military judges within the judicial courts that handle military justice matters.
    \item \textbf{Arbitration Courts:} private entities that resolve litigations outside the public judicial system.
    \item \textbf{Justices of the Peace:} small extrajudicial entities, used to resolve minor disputes.
\end{itemize}


\par Moreover, we have the Public Ministry, an independent entity, representing the State, with the main objective of defending the law and the democratic values of the Portuguese Republic against criminal actors, weather individuals or corporations.

\subsubsection{Sentences, Court Rulings and Jurisprudence in Portugal}
With pertinence to this dissertation, we should also define and distinguish between sentences, \textit{acordãos} (court rulings) and jurisprudence and how are these structured, so we can better apply LLMs to their predictions.

\par A sentence can be defined as a judicial decision made by a singular judge. These are majorly applied in first instance courts and simpler cases. 
On the other hand, \textit{acordãos} or court rulings are decisions made by a collective of three judges, usually in courts of appeal or higher.
So we can define jurisprudence, as the set of past court rulings and sentences, on a specific legal matter, that can be used as supportive material for future rulings.

\par The structure of sentences and \textit{acordãos} is set by the Article 374 of the Decree-Law 78/87 of the Penal Code \parencite{noauthor_lei_78/87}, which states that 
a judicial decision must have three main parts: The \textit{Relatório} (Report), the \textit{Fundamentação} (Foundation/Justification) and the \textit{Dispositivo} (Decision).

\begin{itemize}
    \item \textbf{Relatório:} This section should contain the identification of all parties (defendant, assistants and civil parts) and should serve as a brief contextualization of the case or the appeal if applied.
    \item \textbf{Fundamentação:} Here sits the main body of the case, where the judges must enumerate all facts, proven and not proven, and must clearly expose their motives and reasoning, based on contextual facts, relevant jurisprudence and doctrines and finally the applicable laws, to the given case.
    \item \textbf{Dispositivo:} Finally, they must end, by clearly stating the final decision they have reached, alongside all the legal implications of such decisions (e.g, sentences, fines, compensations, destinations of objects, etc.).
    
\end{itemize}

It's important to note, although this structure is mandatory, which in turn creates certain patterns we can mine, there is still a great degree of freedom, judges have when writing their decisions, creating some variablity and difficulties when mining these documents.
This will be further explored and analyzed in the Methodology chapter.

\par Now that we have a better understanding of the skeletons of sentences and \textit{acordãos}, we have solid foundations on how to apply reasoning techniques to our language models, going beyond just predicting the final decision.
These documents, base themselves on the rigorous legal explainability of the decisions. Such austerity when it comes to justifying a ruling is mandatory, due to the possible negative implications of a judicial error.


% Have a diagram of the portuguese judicial system here. Also a diagram summarizing sentence structure, that is pertinent problably just have that % 









\section{Overview of Large Language Models}
% History, development and recent innovations in LLMs

Over the past years, LLMs have gained substantial interest in various fields, particularly in Law. This inherent relevance derives from the central importance language has when considering human nature, enabling the expression of one's thoughts, emotions, and promoting the communication with others. However, it's important to keep in mind, these models cannot think like humans, rather, emulate language comprehension and generation behind mathematical and statistical frameworks, all to be further explained in the subsequent chapter of Theoretical Foundation. \parencite{wang_history_llms_2024,naveed_comprehensive_overview_llms_2024}.

\par Nevertheless, the chronological historical evolution of Language Models(LMs), date back to Statistical Language Models in  the 1990s \parencite{rosenfeld_two_2000}, where the problem of processing context and semantic dependent features of natural language, was approached from a statistical perspective.


At that time, major barriers were primarily related to the dimensionality curse(i.e., the exponential increase in computational complexity as context size grows) and the disparity between the over-simplistic models (e.g., n-gram model, that predicted a word's likelihood based on previous n-words), and the complex nature of human language. \textcite{rosenfeld_two_2000}, recommended future work to focus on injecting prior human knowledge of language into the models, trough statistical frameworks.

\par Overtime, the resurgence of Neural Networks(NNs), enabled the application of these to the natural language problem. By vectorizing or tokenizing the words and symbols in a given sentence, it was possible to pass these numerical representations through NNs, that mapped similar words to proximate vectors that clustered efficiently semantically similar text, then outputting a probability over the density distribution of words/sentences in a given task. 

\textcite{bengio_neural_2000} had noted that, utilizing a model, where each training phrase, exhibited systematic generalization to unseen word combinations, performing the objective in a more resource efficient manner. They proposed future research to study the implementation of Recurrent Neural Networks (RNNs), to permit taking advantage of the sequential nature of language, which basic NNs could not, and warrant the use of larger training samples (e.g., a whole paragraph, at the time) without increasing excessively the number of parameters of the network.

\par Subsequently, RNN's became the staple when it came to LMs, surpassing its precedents, in terms of accuracy, and the capacity of handling larger amounts of text \parencite{mikolov_recurrent_2010}. Even so RNNs, still struggled, due to the impossibility of parallelization of training examples, inherent to the sequential nature of the network and the  “vanishing” gradient problem in longer context windows (i.e, extreme shrinkage of gradients during training) that reduced learning abilities.

Although architectural improvements were made, namely with the introduction of Long Short-Term Memory and Gated Recurrent Units(GRU), which utilized gate mechanisms to optimize information flow and reduce the aforementioned issues, these networks still struggled with massive context size inputs.

\par A breakthrough, happened when \textcite{vaswani_attention_2023} proposed using the Attention Mechanism (e.g. focusing on specific keywords in a sentence) as the main foundation for identifying patterns between the input and the output, and Self-Attention, capturing dependencies within the input. These architectural innovations, not enabled taking advantage of Graphics Processing Units (GPUs) technology, which are specialized in handling repetitive calculations in a parallel format, reducing computation time. Alongside, they also facilitated the detection of long-range dependencies.
Self-Attention yielded inherently more interpretable models, by showing which input elements were prioritized in a given task. Both of these mechanisms are part of the Transformer Network, that permitted scaling LMs to today's standards. The Transformer and its components will be further explored in the Theoretical Foundation chapter.

\par With the previously mentioned innovation, the scale of training data increased substantially, particularly on unlabeled text (e.g., text found on websites like Reddit, Common Crawl, Wikipedia, etc.).\parencite{wang_history_llms_2024}. Thereby permitted LLMs to grasp fundamental language structure, such as vocabulary, semantics, syntax, and logic of various tongues and dialects. Consequently, LLMs, excelled in various Natural Language Processing(NLP) tasks such as text summarization, question-answering or even translation, all vital for legal applications.

\par More recently, with the emergence of LLMs like GPT-4.5 \parencite{openai2025gpt45} or Gemini \parencite{team_gemini_2025} alongside open source models such as DeepSeek \parencite{deepseek-ai_deepseek-r1_2025} or Llama 4 \parencite{meta2025llama4}, and many others \parencite{shukla_Most_advanced_models_2025}, have pushed the capabilities of what and how these models can be applied. These state-of-the-art models pursue multi-modality, being capable of interconnecting text, images or videos in various tasks and show high reasoning capabilities.

\par Considering the history and recent innovations in the field of LLMs, their use in legal practice seems inevitable, and for the importance of this dissertation we shall explore how have they been applied in Law, particularly in the Portuguese context, and how can they be further explored in this area.



% IMAGE OF THE TIMELINE OF LANGUAGE MODELS

%%%%%%%%%%%%%%%%%%%%%%%%%%%%%%%%%%%%%%%%%%%%%%%%%%%%%%%%%%%%%%%%%%%%%%%%%%%%%%%%%%%%%%%%%%%%%%%%%%%%%%%%%%%%%%%%%%%%%%%%%












\section{LLMs in Legal Application}
% Talk about AI in legal practice annd Reviews statistical/ML approaches and LLM-based prediction globally.
%  IMAGE SHOWING DIFFERENT APPLCIATIONS OF LLMS

AI, particularly LLMs have never been more popular in Law applications. They're use can be segmented into various legal tasks, ranging from document summarization or drafting to more complex tasks such as judicial outcome prediction. This section will explore in more detail recent key applications in the LegalAI field, what are the current limitations and lastly proposed solutions alongside future research directions.

\subsection{Main Use Cases}
The key applications of LLMs in Law practice can be segmented into 3 categories, such as Legal Prediction, Legal Document Review and Legal Research \parencite{siino_Lit_review_LLMs_2025}. We will further explore each main use and their possible branches of applications.

\subsubsection{Legal Prediction}
In this specific task, the literature highlights the natural evolution from using traditional ML methods (e.g. Tree Models, Statistical Models, etc.) to the more advanced use of LLMs in Judgment Prediction(e.g. BERT or GPT based models). This task goal is forecasting the outcome of a legal case, by analyzing a great amount of legal data, ranging from past court decisions, applicable laws and case context to predict a court verdict of current or future diligence's \parencite{siino_Lit_review_LLMs_2025, zhong_summary_LegalAI_2020}.
\par However, Judgment Prediction is a extremely complex problem not only in technical domain, but also in a ethical domain. In a later Section we will discuss further on these dilemmas. 
\par Finally, we can also affirm that this particular problem, is the kernel of my investigation, where I will adapt and possibly innovate on the related work already present in recent literature.

\subsubsection{Legal Document Review}
% re write this, made by ai
This application incorporates a broad range of tasks, which the literature commonly categorizing into either classification, regression or generative problems:
\begin{itemize}
    \item \textbf{Classification Tasks}: 
    \begin{itemize}
        \item \textit{Named Entity Recognition}: Identifying specific entities in text based on context, such as recognizing ``Apple'' as Apple Inc. rather than the fruit in a contract \parencite{siino_Lit_review_LLMs_2025, shao_llms_meet_law_2025}.
        \item \textit{Document Classification}: Categorizing legal documents based on their content, such as identifying contract types (e.g., distinguishing between permanent contract, fixed-term contract, indefinite-term contract, etc.) \parencite{siino_Lit_review_LLMs_2025, shao_llms_meet_law_2025}.

    \end{itemize}

    \item \textbf{Regression Tasks}:
    \begin{itemize}
        \item \textit{Similarity Estimation}:
        The goal is predicting a continuous value (usually between 0 and 1), representing how similar given entities (e.g., documents, court cases, contracts, legal opinions, etc.) are \parencite{siino_Lit_review_LLMs_2025,zhong_summary_LegalAI_2020}.
    \end{itemize}


    \item  \textbf{Generative Tasks}:
    \begin{itemize}

        \item \textit{Document Automation}: Generating or structuring legal documents, such as contracts or motions, improving lawyers and other justice system agents efficiency and scalability, by allowing them to center their attention on more high reasoning legal tasks, instead of mundane document preparation \parencite{siino_Lit_review_LLMs_2025, shao_llms_meet_law_2025}.

        \item \textit{Document Summarization}:
        By shortening the verbosity of a given legal document while retaining the main takeaways, a legal agent can in a more efficient manner analyze such documents \parencite{siino_Lit_review_LLMs_2025, shao_llms_meet_law_2025}.
    \end{itemize}
\end{itemize}

\subsubsection{Legal Research}
Finally, a justice agent, needs to research various legal documents to sustain his arguments. Trough traditional methods like simple Google searches or specific databases searches like \href{https://www.dgsi.pt/}{dgsi.pt} , an experienced agent could get decent work done, but augmenting it's research with LLMs in a responsible manner, the improvements in time efficiency could be beneficial.
\par The main branching tasks at this domain are:

\begin{itemize}
    \item \textbf{Document Retrieval}: Searching for documents in a natural language format, instead of key word/sentences is pivotal in the scalability of an agent work. However this specific task has shown potential misuses in the past with models hallucinating made up court cases or law articles. It's important that document retrieval to be complemented with other techniques like (Retrieval Augmented Generation)RAG or Knowledge Graphs \parencite{siino_Lit_review_LLMs_2025, shao_llms_meet_law_2025}.

    \item \textbf{Case Entailment}: Determining wether a document has verified facts, legal principals and arguments it's crucial to one´s work. Case Entailment complements Document Retrieval by creating a system capable of identifying and associating paragraphs within a specific precedent case, that could be relevant in justifying a argument or decision in a case \parencite{siino_Lit_review_LLMs_2025, shao_llms_meet_law_2025}.

    \item \textbf{Question-Answering}: Lawyers and other agents don't know everything about Law, so it's natural that they want a question answering system. This could be useful not only for a preparation for a case but also for consultancy purposes with clients \parencite{siino_Lit_review_LLMs_2025, shao_llms_meet_law_2025}. Here chat bots like Chat-GPT have been implemented, particularly at firm levels. However, even with the perceived benefits of such system, there are major concerns in it's use like Hallucinations, Data Security and Privacy, Compliance Issues (i.e. can´t give legal advice without professional supervision), among others further explained in \textcite{homoki_llms_use_cases_2024}
    
\end{itemize}


%%%%%%%%%%%%%%%%%%%%%%%%%%%%%%%%%%%%%%%%%%%%%%





\subsection{Current Barriers}
Even though LLMs posses many advantages in Law applications, there are  are both technical and non-technical barriers, currently, that need further research and mitigation.
\par Given the inherent statical nature of LLMs, their outputs could generate occasionally false, yet plausible looking facts or references, such as non existing court cases or law articles \parencite{armour_augmented_lawer_2020,surden_chatgpt_AI_llms_law_2023}, with even Portuguese news reporting on current investigations of this phenomenon in our court system \parencite{publico2025alegado}.Others issues such as limited, bias and novel data, can reduce performance since these models depend on their training data. Also LLMs might be difficult to interpret without certain explainable techniques, which could reduce the trust and accountability of the models outputs. \parencite{armour_augmented_lawer_2020,surden_chatgpt_AI_llms_law_2023}.Moreover, there is the inherent computational cost of running these models, and in dealing with a field, where one could ingest a large volume of data (e.g., a large firm contract), this could result in a steep price to implement \parencite{armour_augmented_lawer_2020,surden_chatgpt_AI_llms_law_2023}.

\par Alongside the technical difficulties presented above, there are other concerns regarding ethical, legal and regulatory domains of LegalAI, such as data privacy and security, bias and discrimination, liability or property rights. These will all explored more, in a section further down.



%%%%%%%%%%%%%%%%%%%%%%%%%%%%%%%%%%%%%%%%%%%%%%%%%%%%%%%%%%%%%%%%%%%%%%%%%

\subsection{Proposed Solutions and Future Directions}
To mitigate some of the technical concerns, the literature unanimously suggests rigorous human supervision and validation of models outputs. At the same time, users must educate themselves on the strengths and limitations of these tools to avoid over reliance and malpractice. Moreover, to tackle the specific problems legal application brings, there must be reinforced research in specializing models in specific law contexts, taking advantage and improving on recent models multi modality and higher reasoning capabilities, complementing LLMs with other techniques like RAG or KG and finally enhancing transparency and interpretability of these models.\parencite{surden_chatgpt_AI_llms_law_2023,siino_exploring_llms_law_2025}

\par Concerning the non-techninal hurdles, there must be investment in secure data handling protocols, contractual safeguards and regulatory frameworks to avoid the pitfalls mentioned. \parencite{surden_chatgpt_AI_llms_law_2023,siino_exploring_llms_law_2025,shao_llms_meet_law_2025} 


\par In my view, the interdisciplinary dialogue between AI experts and legal professionals, will be key, to advance this field.






%%%% To rrefine the above sections #######################################






\subsection{How are LLMs used currently in the Portuguese Judicial Context}
%Go deeper on the portuguese context
In Portugal's judiciary system, there is also a movement towards innovation, as portrayed in the  \href{https://govtech.justica.gov.pt/}{“Estratégia GovTech da Justiça”}, a government initiative to digitally evolve the legal domain in Portugal.
LLM's, due to their capabilities in natural language, as discussed above, become inevitably a vehicle towards that innovation, both in the public and private sector.

\subsubsection{Public Sector}

Recently, there have been a growing number of public entities actively or planning to implement AI solutions in various legal tasks.

\par \textcite{IA_Portugal_2024} discusses the most recent initiatives taken by public institutions regarding the application of LLM's into the justice system, and I will present a few of them. 
\par For example, in 2023, the \textit{Supremo Tribunal de Justiça }(Supreme Court of Justice), presented \href{https://stjiris.github.io/}{"IRIS- Information, Rationalization, Integration and Summarization: Application of Artificial Intelligence Techniques at the Supreme Court of Justice"} , in collaboration with INESC-ID, with multiple objectives, such as assistance in case analysis, preparation, anonymization and finally transcribing of old transcripts to digital format. Alongside this, there was also the goal public dissemination of Supreme Court cases in their website. Per the researchers, this was based in techniques like optical character recognition to extract the text from scanned documents, then complemented with pipelines of document search, retrieval and analysis.

\par Another case, is \href{https://justica.gov.pt/Servicos/Guia-pratico-da-Justica}{Guia Prático de Acesso da Justiça (GPJ)}, a chatbot, as per the public sources, specialized in Marriage and Divorce Litigation and Company Creation laws. It's stated that is based of the GPT-4 model of OpenAI and distributed trough Microsoft Azure cloud platform.

\par Further reading of \textcite{IA_Portugal_2024} is recommend to learn about more recent applications of AI and LLMs in the Portuguese Public Legal System.

\subsubsection{Private Sector}
The portuguese private sector is rapidly evolving when it comes to the application of AI and LLMs, pushed mainly by big law firms, due to the necessity of improving staff efficiency and to address local and global competition in the sector.

\par One of the standout case studies, regards the portuguese startup \href{https://legau.com/}{Legau}, and their collaboration with firms such as CCA Law Firm \parencite{legau_cca_2025}, RFF Lawyers \parencite{legau_rff_2025}, Antas da Cunha ECIJA \parencite{legau_adcecija_2025}, among others. Based on information provided on their product, they leverage LLMs for legal documents research, planning, reviewing, translation and creation, with robust integration with common work tools such as Microsoft Word and Outlook. Antas da Cunha ECIJA and RFF Lawyers reported an increase of 20.7\% and 19\% average time gain per individual working with Legau's tools. 

\par Another growing portuguese LegalAI startup is  \href{https://www.affine.law/pt}{Affine Law by NeuralShift}, with a specialized focus on deep research capabilities of both public legal information and firm specific document database, with claims of 60\% faster research speeds. Some of the partnerships involve firms like Melo Alves, Proença de Carvalho, Vieira de Almeida and Gama Glória \parencite{affine_2025}, with the latter commenting an interesting observation regarding the use of Affine, in my opinion: \textit{" Research is efficient, and the end result is equivalent to having an unlimited number of junior associates dedicated to legal research. This gives the team more freedom to focus on strategic, high-impact tasks, allowing us to achieve more with the same number of lawyers on a project." }
Testimony's like this highlight the real world impact LLM based systems can have when done tailored to customer needs.


\par Meanwhile, Abreu Advogados and Morais Leitão, are firms, that adopt foreign  tools, such a Microsoft 365 Copilot \parencite{abreu_copilot_2025} and Luminance \parencite{mlgts_innovador_2025}, that enable  automated email drafting,  contract processing and analysis, question-answering, among other features, based on the main use cases of LLMs in LegalAI.

\par A common thread between all, is the commitment to data privacy and security, with all giving guarantees that follow EU guidelines of the topic.



















\section{Related Work}
% Go deeper into studies related to mine about judicial outcomes predictions

Regarding the theme of legal judgment prediction, there have been various studies trough time, that have investigated the application of 
of various ML and DL techniques. Unfortunately, I was not able to find any similar study, particularly in the portuguese context, which could signify
the novelty of this dissertation in Portugal.

Nonetheless, based on my research, the most organized way of discussing these findings is 
based on the type of techniques used, from traditional ML methods, to more complex Neural Network architectures, that include LLMs.

\subsection{Classical Machine Learning Approaches}
 \textcite{aletras_predicting_2016} made a foundational study, that has been cited thousands of times, 
where they used primarily a linear kernel Support Vector Machine(SVM) to predict whether there was a violation or not of Article 3 (prohibition of torture), Article 6 (right to a fair trial) and Article 8 (respect for one's privacy and family life) of the European Convention on Human Rights (ECHR).
A N-gram algorithm was used to extract the textual features from the 250 cases regarding Article 3, 80 cases regarding Article 6 and 254 regarding Article 8. 
Finally, they achieved an average accuracy of 79\% across all articles. 
An interesting pattern noted was that the most decisive predictive factors, were based upon a case circumstances, rather than raw legal principles. 
They conclude that this can be an indicative on how judges make decisions, particularly in ethically charged cases, like the one's treated at ECHR.

\par Another study, made in Europe, was \textcite{sulea_predicting_2017}, where they also used linear SVMs to predict a case outcome, but this time on French Supreme Court rulings.
The dataset was made of 126,865 cases,spanning from the 1880s until the 2010s, from various doctrines, such as civil, criminal, social and commercial.
The textual features were extracted using common NLP techniques, such as Bag-of-Words(unigrams) and bi-grams, 
and the results obtained were an accuracy of 96.9\% and a f1-score of 97\%. 
When the researchers compared their results with previous work, they found that their model outperformed various studies, 
but highlighted that true comparisons were difficult, due to language and most importantly jurisdiction differences, where they conclude french justice
leaves less room for interpretation for judges, when compared to the previous studies cited in their paper.
This highlights an key aspect, where the legal context might difficult models generalization, underlining the importance of specialized studies based upon local jurisdiction,
like the one proposed in this thesis.

\par More recently, a study made in Thailand by \textcite{anantathanavit_utilizing_2024}, also employed various classical ML methods 
like Logistic Regression, Decision Trees, Random Forests, SVMs and Naive Bayes, to predict the outcome of Supreme Court Cases. 
This time the textual features were extracted in various ways, such as Count Vectorizer, TF-IDF and Word2Vec and Doc2Vec embeddings. 
These featurizers + ML models were compared to a thai fine-tuned BERT model, and in average the BERT model outperformed all previous combinations, with an accuracy of 74.6\%.
Due to these findings, the researchers recommended the further exploration of DL methods, to handle the complexity of legal documents.

\par These methods, were the first to be applied and now serve as baselines for more complex models approaches. 
The same intuition will be applied in this dissertation, where I will compare the performance of classical ML methods with more complex DL techniques.

\subsection{Neural Network Approaches}
An early study, that applied DL methods to this problem, was \textcite{long_automatic_2018},
used a Bi-directional GRU to encode the textual features, then a pair-wise attention mechanism to model the complex interactions of the
law documents, while increasing its interpretability. Finally, final convolutional layer is employed to predict the final outcome of the case.
Their model was  trained and evaluated on 100,000 from the Supreme People's Court  of China,
and achieved an accuracy of 82.2\% and a f1-score of 83.4\%, surpassing all ML or simpler DL baselines.
An interesting aspect of this study, and many other chinese ones, is their large datasets, commonly taken from the
\href{https://wenshu.court.gov.cn/}{China Judgments Online website}, which countains hundreds of millions of carefully structured legal documents. 
Other examples of massive legal databases are the USA's \href{https://case.law/}{Case Law} or the EU's \href{https://eur-lex.europa.eu/}{EUR-LEX}, which portuguese jurisprudence is also included.
When considering this dissertation data source, \href{https://www.dgsi.pt/}{dgsi.pt}, one must be aware of the smaller size of the database. 
Another interesting field of approach in the portuguese context, could be the enhancement of this public source of legal documents, that could be further explored in future research.

\par Continuing with Neural Network approaches, we can highlight a brazilian study, \textcite{bertalan_using_2024}, where they used various classifiers, 
including a Hiercachical Attention Network(HAN), composed by a mix of GRU and Attention mechanism layers, to predict the outcome of 2000+ homicide and corruption cases in the state of São Paulo.
The HAN model achieved an accuracy of roughly 99\% in both types of cases, but did not significantly outperform other simpler baseline models,
such as a simple Regression Tree, a SVM or a simple GRU network, which achieved accuracies between 95-99\%. One could argue , the deployment of more complex models, like HAN,particularly in a small dataset like this, may not be as 
cost effective, but more research analysis for my hypothesis is needed.
Another brazilian investigation, \textcite{menezes-neto_using_2022}, predicted the reversal or confirmation of appeals in the 5th Regional Federal Court in Brazil. 
This study, had to it's disposal a much larger dataset, than the previous one, with over 700 000 cases. 
They used various DL architectures, such as AWD-LSTM, BERT+LSTM and BigBird, and all surpassed legal experts in terms of Matthews Correlation Coefficient(MCC), with the best model, AWD-LSTM achieveing a MCC of 0.3688 against the experts 0.3531 average MCC score.
Although these values are not distant, by quoting the researchers, we can observe a major justification for these ML/DL models:
\textit{"It would be nearly impossible for our expert panel to evaluate all entries within our validation and test datasets."} Basically, instead of dealing with roughly 76 000 test cases, these legal experts, only dealt with 690.
Ultimately,the relevance of brazilian studies, is due to the lack of similar research in Portugal and also due to the similarities between the portuguese and brazilian legal systems, due to the inheritance of the judicial matrix during the colonization period.

\par Finally we can mention and analyze a few studies regarding the use of LLMs in this problem, since it's the main focus of this thesis.
\par \textcite{kmainasi_can_2025} is a recent study, with very similar objectives and methodology as the proposed in this thesis.
They utilized public data from the Saudi Ministry of Justice, with 3752 training samples and 538 testing samples.
One of the limitations of this investigation is the smaller size dataset, but most importantly the niche language they are in, Arabic, which has less research associated with it, similarly to European portuguese, in this specific problem. 
To achieve their goal of judicial sentence prediction, they compared baseline Llama-3.2-3B and 8B models, only with prompt engineering techniques, like few-shot learning for example, versus the same model with the same instructions, but fine-tuned to the specific legal context.
They found an increase of an BERT-score-F(f1-score based on contextual embeddings to measure semantic similarity) of 54\% and 58\% for the 3B and 8B non fine-tuned models respectively, versus the 74\% and 76\% of the fine-tuned counterparts, concluding that fine-tuning was a key aspect in the perfomance improvement. 
Alongside these quantitative results, they also analyzed the qualitative outputs of the models, and found that the fine-tuned models, were able to justify their predictions better in aspects like legal language, coherence and clarity to name a few, when compared to the non fine-tuned models.
This research, will be a key reference for my dissertation, due to the similarities with my proposed methodology in the LLM modelling.

\par At last I would like to highlight studies that extrapolate the simple use of LLMs, but enhance them with other techniques, such as RAG, to improve models performance and interpretability in case prediction.
Analyzing \textcite{nigam_nyayarag_2025} and \textcite{peng_athena_2024}, both utilized RAG techniques, to improve the use of LLMs, proprietary and open source, respectively.
In the first study,their case text + RAG system with combination of statutes, precedents outperformed the baseline predictions, using only the integral text of the case, with the an accuracy of 65.86\% and a f1-score of 63.96\%. 
In other lexicial and semantical metrics, such as ROUGE-L or BERT-score, the RAG system also outperformed the baseline.
The second study, had a similar approach, and also concluded that their RAG system "Athena", excedeed the baseline LLM model in all evaulation metrics.
These studies, serve as key references, by showing the potential external database systems can have in the task of legal judgement prediction. 
If possible, I will try to implement this technique and compare its predictive but also interpretable power, versus the use of just LLMs.


\par To conclude, with this small review of what's been done in the field of legal outcome prediction,
I hope to have shown the main patterns and methodologies uses. 
They will serve as foundational work to build and adapt to the portuguese legal context.












\section{Ethical and Legal Considerations}
% Addresses bias, fairness, GDPR, and Portuguese legal compliance. .
% Outlines frameworks for responsible AI deployment
\subsection{Concerns}

It's impossible to disscuss aplications of AI, particularly LLMs in Law, without addressing in a thorough manner, the ethical and legal implications that arise from it.

\par Studies like \textcite{remolina_ai_2024} or \textcite{jialin_exploring_2026}, highlight that one of the major risks of using LLMs in law practice, is the perpetuation and amplification of existing biases, inherent in the training data. 
This problem, leads to an increased unfairness in legal outcomes, by enhancing existing prejudices that may exist regarding gender,race,socialeconomic status, etc.

\par Other common concerns, among peers is the lack of transparency and explicability of generative models, due to their complex architectures e processes of decision making, knonwn as the "black box" problem. 
Alongside this, the natural risk of Hallucinations in model outputs, exacerbates the trust issues. 
For example, there have been many reports of misues of LLMs, particularly popular chatbots in the generation of legal analysis, outputting inexistent laws, jurisprudence or doctrine. 
Even in portugal, there have been news reporting such cases like \textcite{publico2025alegado}. In my opinion, this problem goes beyond model technical issues, and bleeds into user-side malpractice, usually associated with lack of literacy regarding these tools.
In my study, I will try to mitigate these issues, by using explainability techniques and if possible validation by legal experts, to evauluate if the outputs are coherent, justifiable and trustworthy.

\par The last cited concern above, leads to the next major ethical question: who is held accountable for artifically generated content, particularly in the legal practice? For example, if a judge uses an LLM or any AI tool, to analyze, draft or even decide over a given case, is the judge liable? Could it be the developers or the suppliers of the tool? 

\subsection{How to combat this}

% flr de portugal, eu, ai act, gdpr, etc

To mitigate these concerns...

%minha opiniao baseada na literatura em relação a este tema
In the justice system, inconsistencies in decisions are not easily tolerated,and considering the inherent probabilistic nature of LLMs, for me is obvious that these tools, cannot be used in an fully autonomous faction, without human supervision, for example by a judge or a lawyer.
Alongside this, there needs to be a clarification regarding a common myth, wrongly propagated at times, where tools like LLMs, are seen as replacements for legal professionals.





\section{Gaps in the Literature}
% Highlights the lack of studies on Portuguese law and error detection with LLMs.

\section{Conclusion of Literature Review}
% Summarizes findings and positions this research.



%%%%%%%%%%%%%%%%%%%%%%%%%%%%%%%%%%%%%%%%%%%%%%%%%%%%%%%%%%%%%%%%%%%%%%%%%%%%%%%%%%%%%%%%%%%%%%%%%%%%%%%%%%%%


% Chapter 3: Theoretical Framework
\chapter{Theoretical Foundation}

\section{Introduction to Theoretical Framework}
% Presents foundational theories and concepts.

\section{Natural Language Processing (NLP) and LLMs}
% Explains Statistical Language Models, Neural Netwokrks, RNN's and  Transformer architecture and key models.


% \section{Predictive Modeling}
% % Covers classification techniques and metrics (e.g., accuracy, F1-score).
% This is too simple to be worth a section

\section{Recent Techniques to improve LLM performance}
% Slight intro that zero shot prompting at more complex tasks such as legal may not be enough, so introduce other techniques 

\subsection{Prompt Engineering}
% Introduce what it is and its recent developments, particularly in law tasks

\subsubsection{Few Shot Learning}
% Few shot learning explained
\subsubsection{Chain-of-Thought (CoT) Prompting}
% Describes structured reasoning for judicial analysis.
\subsubsection{Talk about other relevant Prompt engineering techniques to improve model perf....}


% \section{Retrieval-Augmented Generation (RAG)}
% % Details RAG’s role in enhancing legal context retrieval.

% \section{Knowledge Graphs}
% % Explain what it is and how it can be used alongside RAG to improve models performance in law

% \section{Fine-tuning}
% % State of the art fine tuning techniques and its pros and cons


\section{Explainable AI (XAI)}
% Introduces tools (e.g., SHAP, LIME, Attention weights) for interpretable outputs in both NLP and LLMs.






% Chapter 4: Methodology
\chapter{Methodology}
\section{Introduction to Methodology}
% Describes the research design and approach.


\section{Data Preparation}
\subsection{Data Selection}
% Uses public rulings from dgsi.pt; explores feasibility of a manual RAG database of laws.
\subsection{Data Extraction}
% Web scraping and API usage for data collection.
\subsection{Data Preprocessing}
% Cleans, anonymizes, and structures legal texts.


\section{Models}
\subsection{Baseline Models}
% Traditional NLP&ML techniques with basic out of the box encoders for benchmarking.
\subsection{LLM Modeling}
% Chooses lightweight models for 8GB VRAM GPUs (e.g., DistilBERT, TinyLLaMA).

\section{Model Improvements}
\subsection{Fine-Tuning Process}
Adapts models for Portuguese legal texts.
\section{Predictive Modeling}
\subsection{Model Architecture}
Integrates RAG for context-aware predictions.


\subsection{Training and Validation}
Splits data into training/validation/test sets.

\subsection{CoT for Error Reasoning}
Uses CoT to explain detected errors.
\section{Explainability and Validation}
\subsection{XAI Techniques}
Applies SHAP/LIME for model transparency.
\subsection{Lawyer Validation}
Collaborates with law firm for case filtering and performance validation.
\section{Evaluation Metrics}
Includes precision, recall, F1-score for prediction and error detection.
\section{Ethical Considerations}
Mitigates bias and ensures GDPR compliance.

% Chapter 5: Experiments and Results
\chapter{Experiments and Results}
\section{Introduction to Experiments}
Describes experimental setup and goals.
\section{Predictive Modeling Results}
\subsection{Performance Metrics}
Reports accuracy, F1-score across domains.
\subsection{Comparative Analysis}
Compares with baselines (e.g., non-LLM models).
\section{Error Detection Results}
\subsection{Detection Accuracy}
Reports precision, recall, F1-score for errors.
\subsection{Case Studies}
Analyzes specific error detection examples.
\section{Explainability Analysis}
\subsection{XAI Outputs}
Visualizes model explanations.
\subsection{Lawyer Feedback}
Summarizes law firm’s qualitative input.
\section{Discussion of Results}
Interprets findings and addresses limitations.

% Chapter 6: Discussion
\chapter{Discussion}
\section{Introduction to Discussion}
Contextualizes results within legal tech and AI.
\section{Implications for the Portuguese Judicial System}
Explores potential to enhance efficiency and accuracy.
\section{Contributions to Legal Tech and AI}
Details advances in LLM applications for civil law.
\section{Ethical and Legal Implications}
Reflects on bias and compliance challenges.
\section{Limitations and Challenges}
Discusses data, GPU, and law firm constraints.
\section{Future Research Directions}
Proposes expanding domains, improving XAI, refining RAG.

% Chapter 7: Conclusion and Future Work
\chapter{Conclusion and Future Work}
\section{Summary of Findings}
Recaps key results and insights.
\section{Academic Contributions}
Highlights novelty in legal AI for Portugal.
\section{Practical Contributions}
Emphasizes potential to improve judicial processes.
\section{Final Reflections}
Shares lessons learned during the research process.
\section{Future Work}
Suggests next steps: refine models, expand datasets, explore practical implementations.

% Adding bibliography
\backmatter
\printbibliography[title={References}]

% Including appendices
\begin{appendices}
\chapter{Data Samples}
Sample rulings from dgsi.pt.
\chapter{Code and Models}
Links to code repositories or snippets.
\chapter{Lawyer Validation Reports}
Feedback summaries from law firm collaboration.
\end{appendices}

\end{document}
